% Figure: the definiteness map of a 2x2 symmetric matrix in the
% (trace, determinant) plane, with the parabola tr^2 = 4 det marking
% the real-repeated-eigenvalue boundary and the four sign regions.
\begin{figure}[htbp]
\centering
\begin{tikzpicture}
\begin{axis}[
  width=13cm,
  height=8.5cm,
  xlabel={trace $\tau = \lambda_1 + \lambda_2$},
  ylabel={determinant $\delta = \lambda_1\lambda_2$},
  xmin=-5, xmax=5,
  ymin=-4, ymax=6,
  xtick={-4,-2,0,2,4},
  ytick={-4,-2,0,2,4,6},
  tick label style={font=\footnotesize},
  label style={font=\small},
  title={Definiteness of a symmetric $2\times2$ matrix in the $(\tau,\delta)$ plane},
  title style={font=\bfseries},
  axis lines=middle,
  ymajorgrids=true,
  grid style={gray!20},
]
% The parabola delta = tau^2/4 is the repeated-eigenvalue boundary.
% Above it eigenvalues are complex (not for symmetric), on/below real.
% For symmetric matrices eigenvalues are always real, so delta <= tau^2/4.
\addplot[engamber, thick, domain=-5:5, samples=80] {x^2/4};
\node[font=\scriptsize, engamber, anchor=south west, rotate=52]
  at (axis cs:3.1,2.4) {$\delta=\tau^2/4$ (repeated $\lambda$)};
% delta = 0 axis (already an axis line). Shade regions with node labels.
% Positive definite: delta > 0 and tau > 0.
\node[font=\footnotesize\bfseries, engblue, align=center]
  at (axis cs:3.2,3.6) {positive\\ definite};
\node[font=\scriptsize, engblue] at (axis cs:3.2,2.9)
  {$\lambda_1,\lambda_2>0$};
% Negative definite: delta > 0 and tau < 0.
\node[font=\footnotesize\bfseries, engred, align=center]
  at (axis cs:-3.2,3.6) {negative\\ definite};
\node[font=\scriptsize, engred] at (axis cs:-3.2,2.9)
  {$\lambda_1,\lambda_2<0$};
% Indefinite: delta < 0 (any tau).
\node[font=\footnotesize\bfseries, align=center]
  at (axis cs:0,-2.2) {indefinite\\ $\lambda_1>0>\lambda_2$};
% Semidefinite boundary: delta = 0.
\node[font=\scriptsize, anchor=west] at (axis cs:0.2,0.35)
  {$\delta=0$: semidefinite};
\end{axis}
\end{tikzpicture}
\caption[Definiteness map in the trace-determinant plane]{The
definiteness of a symmetric $2\times2$ matrix read entirely from its
two coordinate-free invariants, the trace $\tau = \lambda_1+\lambda_2$
and the determinant $\delta = \lambda_1\lambda_2$. Sylvester's
criterion becomes a picture: the determinant sign splits definite from
indefinite, and among the definite matrices the trace sign splits
positive from negative. The amber parabola $\delta = \tau^2/4$ is the
repeated-eigenvalue boundary; a symmetric matrix always sits on or
below it, since its eigenvalues are real. A matrix in the upper-right
wedge (both invariants positive) is positive definite, its energy a
bowl; crossing the horizontal axis into $\delta < 0$ turns it
indefinite, its energy a saddle, the flip that the worked example
produces by changing a single diagonal entry. The map is the
two-dimensional face of the general eigenvalue test for definiteness.}
\label{fig:vol02:ch10:definiteness-map}
\end{figure}
